\section{Conclusion}
\label{sec:conclusion}

The single picture at this point is this. There is one substance, experience, in the
form of vibes. The least arbitrary way to describe its structure begins with the whole numbers,
whose symmetries, generalized to reflections, give a hyperbolic crystal whose dimension is
narrowed to a small window and then settled at three. On that crystal each vibe carries a
three-valued tone, relates to its neighbors by noting them, and is updated by one local rule,
with no continuous weights and no hidden state, and with tone, note, fill, will, and link all
being the one vibe seen from different angles rather than added parts. The mesh grows forever at
its present edge, deterministically and without randomness, yet it is free in the only sense
that matters, self-determining and irreducible. From this, space and time and matter and force
and gravity and the quantum and the cosmos emerge as patterns, the variety of experience emerges
as structure on the one tone, and minds emerge as the recursion of the same form, vibes
gathering into higher vibes, selves made of selves, in a tower that runs from a single quantum of
feeling up to the whole, with the world able to model itself only in summary and so without
paradox.

The aim has been completeness of the conceptual picture, not finality of the science. The
physical, the mental, and the spiritual are, on this view, not three substances but one,
appearing in many regions and in two conditions, the lawful that has settled into stable pattern
and the wild that stays fluid, with the physical world being the lawful, shared region we hold in
common and the inner and the unseen being other regions of the same field, lawful in places and
wild in others. We have tried to be clear about what is solid, what is partial, and what is open,
and to keep the calm distinction between describing the structure of feeling and possessing the
feeling itself.

If there is a single thing to carry away, it is the reversal at the root. The usual order puts
dead matter first and hopes mind appears. This picture puts experience first and shows, as far as
a finite model can, that matter is what experience looks like when it settles into stable, shared
arrangements. Whether that reversal is true of our universe is a question for evidence and for
the harder modeling still ahead. That it can be made into a definite, buildable, checkable
picture, rather than a mood, is what this paper has tried to show.
